\chapter{4x10\textsuperscript{9}D.C.}\label{4x10^9-3}

\hfill\break

Stavamo viaggiando da Padang verso l'interno - almeno così avevo capito
- in salita, su strade che a volte erano lisce come la seta e a volte
irregolari e piene di buche, finché la macchina si fermò davanti a
qualcosa che al buio sembrava un bunker di cemento ma, a giudicare dalla
mezzaluna dipinta di rosso sotto una lampada abbagliante al tungsteno,
doveva essere una specie di clinica medica. L'autista, vedendo la nostra
effettiva destinazione, si contrariò - era un'ulteriore prova del fatto
che fossi malato e non semplicemente ubriaco - ma Diane gli mise in mano
altre banconote e lo liquidò, lasciandolo rabbonito se non proprio
felice.

Avevo difficoltà a stare in piedi. Mi appoggiai a Diane, che si accollò
coraggiosamente il mio peso, e rimanemmo in mezzo alla notte umida, su
quella strada deserta, con la luce lunare che penetrava attraverso i
brandelli di nuvole. Di fronte a noi c'era la clinica, un benzinaio
dall'altra parte del marciapiede, e nient'altro che foreste e terreni
pianeggianti che dovevano essere campi coltivati. Non registrammo
presenze umane fino a quando la zanzariera della porta d'ingresso della
clinica si aprì cigolando e una donna bassa e paffuta che indossava una
gonna lunga e un cappellino bianco ci venne incontro in fretta.

«Ibu Diane!» disse la donna, con tono emozionato ma delicato, come se
temesse di essere sentita persino a quell'ora tarda.

«Benvenuti!»

«Ibu Ina» disse Diane con fare rispettoso.

«E questo dev'essere\ldots?»

«Pak Tyler Dupree. Quello di cui ti ho parlato.»

«Sta troppo male per riuscire a parlare?»

«Sta troppo male per dire qualcosa di sensato.»

«Allora portiamolo dentro senza perdere altro tempo.»

Diane mi sorreggeva da un lato e la donna che aveva chiamato Ibu Ina mi
afferrò il braccio destro dalla spalla. Non era giovane, ma aveva molta
forza. I capelli sotto il cappello bianco erano grigi e diradati.
Profumava di cannella. A giudicare da come storceva il naso, io avevo un
odore peggiore.

Quindi entrammo. Passammo per una sala d'attesa arredata solo con mobili
in rattan e dozzinali sedie di metallo, e ci ritrovammo in un
ambulatorio piuttosto moderno, dove Diane mi scaricò su un tavolo
imbottito. Ina disse: «Bene, allora, vediamo cosa possiamo fare per lui»
e mi sentii abbastanza al sicuro da svenire.

\separatorefiori

Mi svegliai al suono di una chiamata alla preghiera proveniente da una
moschea lontana, immerso nel profumo di caffè appena fatto.

Giacevo nudo su un pagliericcio in una piccola stanza di cemento con una
finestra che era l'unico punto da cui filtrava luce, una luce fioca,
pallida premonizione dell'alba. C'era un'entrata coperta da una specie
di passamaneria di bambù e da lì dietro proveniva il rumore di qualcuno
che sbatacchiava energicamente tazze e scodelle.

Gli abiti che indossavo la sera prima erano stati lavati e li ritrovai
piegati vicino al pagliericcio. Ero nello stato di lucidità tra
un'ondata di febbre e l'altra - avevo imparato a distinguere queste
piccole oasi di benessere - e avevo forza a sufficienza per riuscire a
vestirmi.

Stavo in equilibrio su una gamba e prendevo la mira dei pantaloni con
l'altra quando Ibu Ina sbucò dalla tenda. «Allora stai abbastanza bene
da alzarti!» disse.

Durò poco. Caddi all'indietro sul pagliericcio, mezzo vestito. Ina entrò
nella stanza con un vassoio di legno su cui erano poggiati una ciotola
di riso bianco, un cucchiaio e una tazza di metallo smaltato.
S'inginocchiò accanto a me, come chiedendomi se ne volessi un po'.

Mi resi conto di volerne. Per la prima volta da molti giorni ero
affamato. Forse era un buon segno. I pantaloni mi stavano tremendamente
larghi, le costole erano sporgenti da fare impressione. «Grazie» dissi.

«Ci hanno presentato ieri sera,» disse, allungandomi la ciotola, «ti
ricordi? Mi scuso per la sistemazione rozza. Questa stanza serve a
nascondere, non a ospitare.»

Doveva avere tra i cinquanta e i sessant'anni. Il viso era rotondo e
rugoso, i lineamenti si concentravano in una luna di carne bruna,
sembrava una bambola fatta di mele secche e questo aspetto era
accentuato dal lungo abito nero e dal cappello bianco. Se gli Amish si
fossero stabiliti nella parte occidentale di Sumatra avrebbero prodotto
qualcosa di simile a Ibu Ina.

Il suo accento aveva la cadenza indonesiana ma la sua dizione era
perfettamente corretta. «Parli molto bene» dissi, fu l'unico complimento
che mi venne in mente con così poco preavviso.

«Grazie. Ho studiato a Cambridge.»

«Inglese?»

«Medicina.»

Il riso era insipido ma buono. Feci del mio meglio per finirlo.

«Magari ne vuoi ancora, più tardi?» disse Ibu Ina.

«Sì, grazie.»

\emph{Ibu} era un termine minang usato come forma di rispetto nel
rivolgersi alle donne. (L'equivalente maschile era \emph{Pak}.) Il che
implicava che Ina era un medico Minangkabau e che ci trovavamo nelle
regioni montuose di Sumatra, probabilmente nelle immediate vicinanze del
monte Merapi. Tutto quello che sapevo della gente di Ina lo avevo
imparato da una guida di Sumatra che avevo letto sul volo da Singapore:
c'erano più di cinque milioni di Minangkabau che vivevano nei paesi e
nelle città montuose; molti dei migliori ristoranti di Padang erano
gestiti dai Minangkabau; erano famosi per la cultura matrilineare, la
scaltrezza negli affari e l'amalgama tra la fede islamica e i costumi
tradizionali \emph{adat}.

Nessuna di queste cose spiegava però cosa ci facessi nella stanza sul
retro dello studio di un medico Minangkabau.

Dissi: «Diane sta ancora dormendo? Perché non capisco\ldots»

«Ibu Diane ha preso l'autobus per tornare a Padang, mi dispiace. Ma qui
sarai al sicuro.»

«Speravo che anche lei fosse al sicuro.»

«Sarebbe certamente più al sicuro qui che in città, non c'è dubbio. Ma
questo non vi porterebbe fuori dall'Indonesia.»

«Come hai conosciuto Diane?»

Ina sorrise. «Un caso fortuito! O più che altro fortunato. Stava
negoziando un contratto con il mio ex marito, Jala, che tra le altre
cose si occupa di import-export, quando mi resi conto che la nuova
Reformasi era troppo interessata a lei. Per alcuni giorni al mese lavoro
nell'ospedale di Stato di Padang, e fui molto felice quando Jala mi
presentò Diane, anche se stava solo cercando un posto dove nascondere in
via provvisoria una potenziale cliente. È stato così bello conoscere la
sorella di Pak Jason Lawton!»

Era tutto sbalorditivo da quasi ogni punto di vista. «Conosci Jason?»

«Conosco delle cose su di lui\ldots{} A differenza tua, non ho mai avuto
il privilegio di parlare con lui. Ma seguivo attentamente tutte le
notizie su Jason Lawton sin dagli inizi dello Spin. E tu eri il suo
medico personale! E ora ti trovi nella stanza sul retro della mia
clinica!»

«Non credo che Diane abbia fatto bene a darti quelle informazioni.» Ero
certo che non avrebbe dovuto. La nostra unica forma di protezione era
l'anonimato e ora era compromesso.

Ibu Ina sembrò delusa. «Certo,» disse, «sarebbe stato meglio non fare
\emph{quel} nome. Ma gli stranieri con problemi legali sono davvero
all'ordine del giorno a Padang. Esiste un'espressione: ce ne sono a
palate. La situazione degli stranieri con guai legali e di salute invece
è più problematica. Diane deve aver saputo che io e Jala eravamo due
grandi ammiratori di Jason Lawton\ldots{} potrebbe essere stato solo un
atto di disperazione quello di fare il suo nome e comunque non le
credetti fino in fondo, almeno finché non cominciai a cercare delle foto
in Internet. Suppongo che uno degli svantaggi della celebrità debba
essere questa costante esposizione mediatica nelle fotografie. Ad ogni
modo, c'era una foto della famiglia Lawton fatta all'inizio dello Spin
ma la riconobbi: era vero! E allora doveva essere vero anche tutto
quello che mi aveva raccontato del suo amico malato. Eri il medico di
Jason Lawton e ovviamente dell'altro, il più famoso\ldots»

«Sì.»

«Il piccolo uomo nero rugoso.»

«Sì.»

«Il cui farmaco ti fa stare male.»

«Il cui farmaco, almeno spero, mi farà anche stare meglio.»

«Come ha fatto con Diane, a detta sua. Tutto questo mi interessa. Esiste
davvero un'età adulta oltre l'età adulta? Come ti senti?»

«Francamente potrei stare meglio.»

«Ma il processo non è finito.»

«No. Il processo non è finito.»

«Quindi devi riposarti. C'è qualcos'altro che ti posso portare?»

«Avevo dei quaderni\ldots{} dei fogli\ldots»

«In un pacchetto assieme agli altri bagagli. Te li porto. Sei uno
scrittore oltre che medico?»

«Solo temporaneamente. Devo mettere dei pensieri su carta.»

«Forse, quando ti sentirai meglio, potrai condividere con me alcuni di
quei pensieri.»

«Forse sì. Sarebbe un onore.»

Si alzò dalla posizione inginocchiata. «Specialmente sul piccolo uomo
nero rugoso. L'uomo che veniva da Marte.»

\separatorefiori

Dormii in modo irregolare per un paio dei giorni che seguirono, sorpreso
dallo scorrere del tempo, le notti improvvise e le mattine inaspettate,
identificando ciò che potevo degli orari grazie alla chiamata alla
preghiera, al rumore del traffico, a Ibu Ina che mi offriva riso, uova
al curry e spugnature periodiche. Parlavamo ma le conversazioni
scivolavano via dalla mia memoria come sabbia attraverso un setaccio, e
riuscivo a capire dalla sua espressione che di tanto in tanto mi
ripetevo o mi dimenticavo ciò che mi aveva detto. Luce e buio, luce e
buio; poi, all'improvviso, Diane era in ginocchio accanto a Ina vicino
al letto ed entrambe mi lanciavano sguardi tristi.

«È sveglio» disse Ibu Ina. «Vi prego di scusarmi. Vi lascio soli.»

Quindi accanto a me rimase solo Diane.

Indossava una camicetta bianca, una sciarpa dello stesso colore sui
capelli scuri e pantaloni blu svolazzanti. Sarebbe potuta sembrare una
delle normali frequentatrici laiche del centro commerciale di Padang,
sebbene fosse troppo alta e pallida per riuscire a ingannare qualcuno.

«Tyler,» disse, i suoi occhi erano blu e spalancati, «stai facendo
attenzione ai fluidi?»

«Ho un aspetto così brutto?»

Mi accarezzò la fronte. «Non è facile, vero?»

«Non mi aspettavo che fosse indolore.»

«Un altro paio di settimane e sarà tutto finito. Fino ad allora\ldots»

Non c'era bisogno che me lo dicesse. Il farmaco stava cominciando a
lavorare nelle profondità del tessuto muscolare, in quello nervoso.

«Ma questo è un bel posto dove stare,» aggiunse, «abbiamo antispastici e
analgesici decenti. Ina capisce cosa sta succedendo.» Sorrise in modo
triste. «Tuttavia\ldots{} non è proprio come avevamo pianificato.»

Avevamo puntato sull'anonimato. Per un americano danaroso, qualunque
città tra quelle del porto dell'Arco avrebbe dovuto essere un luogo
sicuro in cui perdersi. Ci eravamo sistemati a Padang non solo per la
sua convenienza - Sumatra era la massa continentale più vicina all'Arco
- ma perché la sua crescita economica iperveloce e i recenti problemi
con il governo della nuova Reformasi a Jakarta avevano trasformato la
città in un'anarchia efficiente. Avrei sofferto la fase di decorso del
farmaco in qualche albergo anonimo, e quando fosse finita - quando mi
fossi effettivamente ricostituito - avremmo pagato il passaggio verso un
luogo in cui non ci sarebbe potuto accadere niente di male. Era così che
sarebbero dovute andare le cose.

Ciò che non avevamo considerato era la sete di vendetta
dell'amministrazione Chaykin e la sua determinazione a fare di noi un
esempio - sia per i segreti che avevamo tenuto per noi che per quelli
che avevamo già divulgato.

«Suppongo di essermi resa un po' troppo visibile nei posti sbagliati»
disse Diane. «Avevo registrato noi due con due diversi collettivi
\emph{rantau}, ma entrambi gli accordi sono saltati, improvvisamente la
gente ha iniziato a non rivolgermi più la parola; stavamo attirando
troppa attenzione. Il consolato, la nuova Reformasi e la polizia locale
avevano la nostra descrizione. Non era del tutto \emph{accurata}, ma
abbastanza vicina alla realtà.»

«Ecco perché hai detto a queste persone chi siamo.»

«L'ho detto perché già lo sospettavano. Non Ibu Ina ma di sicuro Jala,
il suo ex. Jala è un tipo molto astuto. Gestisce una società di
spedizioni abbastanza rispettabile. Gran parte del cemento e dell'olio
di palma che transitano per il porto di Teluk Bayur passa anche per
qualche magazzino di Jala. Il settore del \emph{rantau gedang} raccoglie
meno soldi ma è esentasse, e tutte quelle navi piene d'immigrati non
tornano vuote. Arrotonda alla svelta col mercato nero dei bovini e delle
capre.»

«Sembra un uomo che sarebbe ben felice di venderci alla nuova
Reformasi.»

«Ma noi lo paghiamo meglio. E presentiamo meno problemi legali, sempre
che non ci catturino.»

«Ina approva tutto questo?»

«Approva cosa? Il \emph{rantau gedang}? Ha due figli e una figlia nel
Nuovo Mondo. O Jala? Lo ritiene più o meno affidabile\ldots{} se lo
paghi non ti tradisce. O noi? Pensa che siamo a un passo dalla santità.»

«A causa di Wun Ngo Wen?»

«Sostanzialmente, sì.»

«Sei stata fortunata a incontrarla.»

«Non è solo fortuna.»

«Comunque dobbiamo andarcene al più presto possibile.»

«Appena starai meglio. Jala ha una nave pronta a salpare. La
\emph{Capetown Maru}. Ecco perché sto facendo la spola tra qui e Padang.
Ci sono altre persone che devo pagare.»

Ci eravamo trasformati rapidamente da stranieri con i soldi a stranieri
che una volta avevano i soldi. «Tuttavia,» dissi, «vorrei\ldots»

«Cosa?» Mi passò un dito sulla fronte, avanti e indietro, languidamente.

«Vorrei non dover dormire da solo.»

Fece una risatina e mi poggiò la mano sul petto. Sulla mia gabbia
toracica emaciata, sulla mia pelle ancora ruvida come quella di un
coccodrillo e orribile. Non esattamente un invito all'intimità. «Fa
troppo caldo per coccolarsi.»

«Troppo caldo?»

Stavo tremando.

«Povero Tyler» disse.

Volevo dirle di fare attenzione. Ma chiusi gli occhi e quando li aprii
se n'era andata un'altra volta.

\separatorefiori

Come prevedibile, il peggio doveva ancora venire, anche se nei giorni
successivi mi sentii molto meglio: ero nell'occhio del ciclone, così lo
aveva chiamato Diane. Era come se il farmaco marziano e il mio corpo
avessero concordato una tregua temporanea ed entrambi stessero serrando
le fila per la battaglia finale. Cercai di approfittarne.

Mangiai tutto ciò che mi offriva Ina e di tanto in tanto camminavo per
la stanza provando a canalizzare un po' di forza nelle mie gambe ossute.
Se fossi stato in forze, questa scatola di cemento (che era stato il
deposito delle attrezzature mediche di Ina, prima che installasse un
sistema d'allarme nella sua clinica) mi sarebbe sembrata la cella di una
prigione. Date le circostanze, risultava quasi accogliente. Impilai in
un angolo le nostre valigie rigide e le usai come una sorta di
scrivania, seduto su un tappetino di canne mentre scrivevo. La finestra
alta faceva entrare uno spicchio di luce.

Faceva passare anche la faccia di uno studentello del posto che avevo
beccato un paio di volte a sbirciare. Quando lo accennai a Ibu Ina
annuì, sparì per un paio di minuti e riapparve col ragazzo al seguito:
«Questo è En» disse quasi lanciandolo verso di me attraverso la tenda.
«En ha dieci anni. È molto intelligente. Vuole diventare medico, un
giorno. È anche il figlio di mio nipote. Sfortunatamente ha la
maledizione della curiosità a scapito della sensibilità. Si è
arrampicato sui bidoni della spazzatura per vedere cosa nascondevo nella
stanza sul retro. Imperdonabile. Chiedi scusa al mio ospite, En.»

En abbassò tanto la testa che temetti gli scivolassero dal naso gli
enormi occhialoni. Farfugliò qualcosa.

«In inglese» disse Ina.

«\emph{Scusa}!»

«Non elegante ma dritto al punto. Forse En può fare qualcosa per te, Pak
Tyler, per farsi perdonare per il cattivo comportamento?»

En era chiaramente nei guai. Cercai di aiutarlo a uscirne. «A parte
rispettare la mia privacy, niente.»

«Rispetterà certamente la tua privacy d'ora in avanti\ldots{}
\emph{vero}, En?» En si fece piccolo e annuì. «A ogni modo, ho un
lavoretto per lui. En viene alla clinica quasi tutti i giorni. Se non
sono troppo occupata, gli mostro un po' di cose. L'atlante d'anatomia
umana. La cartina al tornasole che cambia colore nell'aceto. En dichiara
di essere grato di questi piaceri.» En annuiva in maniera quasi
spasmodica. «In cambio, quindi, come forma di compensazione per la sua
rozza negligenza da comune \emph{budi}, En diventerà la vedetta della
clinica. En, sai cosa significa?»

En smise di annuire e sembrò diffidente.

«Significa,» disse Ibu Ina, «che d'ora in avanti userai la tua
attenzione e curiosità per una buona causa. Se qualcuno viene nel
villaggio a fare domande sulla clinica - chiunque provenga dalla città,
intendo, soprattutto se sembrano o si comportano da poliziotti - dovrai
\emph{immediatamente} correre qui e riferirmelo.»

«Anche se sono a scuola?»

«Dubito che la nuova Reformasi ti disturberà a scuola. Quando sei in
classe, stai attento alle lezioni. In tutti gli altri momenti, che sia
per strada o a un \emph{warung}, se vedi qualcosa o senti parlare di me
o della clinica o di Pak Tyler (che tu \emph{non} devi mai nominare),
vieni subito alla clinica. Capito?»

«Sì» rispose En, e mormorò qualcos'altro che non riuscii a sentire.

«No,» ribatté prontamente Ina, «non è previsto alcun tipo di pagamento,
che domanda scandalosa! Però, se sarò soddisfatta, potrei farti dei
regalini. In questo momento non sono per niente soddisfatta.»

En se la svignò, la sua maglietta bianca e troppo grande gli si
rigonfiava sulla schiena.

Prima che facesse buio cominciò a piovere, una forte pioggia tropicale
che durò per giorni, durante i quali scrissi, dormii, mangiai, camminai,
resistetti.

\separatorefiori

Ibu Ina mi fece le spugnature durante una notte buia di pioggia,
strofinando via una muta di pelle morta.

«Raccontami che cosa ti ricordi di loro» disse. «Raccontami cosa vuol
dire crescere insieme a Diane e Jason Lawton.»

Ci pensai su. O, per meglio dire, mi immersi nel sempre più torbido
laghetto della memoria a cercare qualcosa da offrirle, qualcosa di vero
ed emblematico allo stesso tempo. Non riuscii a pescare esattamente ciò
che volevo ma qualcosa affiorò in superficie: un cielo stellato, un
albero. Quell'albero era un pioppo argentato, oscuro e misterioso. «Una
volta siamo andati in campeggio» dissi. «Prima dello Spin, ma non
tanto.»

Era bello sentire che la pelle morta veniva lavata via, almeno
all'inizio, ma il derma che emergeva era sensibile, crudo. La prima
passata della spugna era confortante, la seconda sembrava iodio su un
taglio fatto con la carta. Ina se ne rendeva conto.

«Solo voi tre? Non eravate troppo giovani per andare in campeggio? Nel
posto da cui venite vi lasciano tanta libertà fin da ragazzini? O
viaggiavate con i vostri genitori?»

«Non con i nostri genitori. E.D. e Carol andavano in vacanza una volta
l'anno, villaggi turistici o crociere, preferibilmente senza figli.»

«E tua madre?»

«Preferiva stare a casa. Fu una coppia che abitava in fondo alla strada
a portarci negli Adirondacks insieme ai loro due figli maschi, ragazzini
che non volevano avere niente a che fare con noi.»

«Allora perché\ldots{} ah, suppongo che il padre volesse ingraziarsi
E.D. Lawton? Magari chiedere un favore?»

«Qualcosa del genere. Non ho mai chiesto. Nemmeno Jason. Diane forse lo
sapeva\ldots{} lei prestava attenzione a certe cose.»

«Non è poi così importante. Siete andati in un campeggio montano? Rotola
sul fianco, per favore.»

«Il tipo di camping con il parcheggio. Non esattamente la natura
incontaminata. Ma era un fine settimana di settembre e quel posto era
quasi tutto a nostra disposizione. Abbiamo montato le tende e acceso il
fuoco. Gli adulti\ldots» Mi tornò in mente il nome. «I Fitch cantavano
delle canzoni e ci fecero partecipare ai cori. Dovevano avere bei
ricordi dei campi estivi. In realtà era alquanto deprimente. I giovani
Fitch odiavano tutta la situazione e si rintanarono nelle tende con la
musica nelle cuffie. I Fitch più vecchi alla fine si arresero e andarono
a dormire.»

«E vi lasciarono tutti e tre attorno al falò morente. Era una notte
limpida o piovosa, come questa?»

«Una notte limpida d'inizio autunno.» Non aveva niente a che vedere con
quella, con i suoi cori di rane e le gocce di pioggia come proiettili
sul tetto sottile. «Niente Luna ma tantissime stelle. Non calda ma non
troppo fredda, anche se eravamo in un certo senso in alto sulle colline.
Ventosa. Abbastanza ventosa da poter sentire gli alberi che parlavano
tra loro.»

Il sorriso di Ina si spalancò. «Gli alberi che parlano tra loro! Sì,
conosco quel suono. Ora sul fianco sinistro, per favore.»

«Il viaggio era stato noioso, ma cominciava a piacerci adesso che
eravamo solo noi tre. Jase afferrò una torcia elettrica e ci
allontanammo dal fuoco di qualche metro, andando verso uno spazio aperto
in un bosco di pioppi, lontani dalle auto e dalle tende e dalla gente,
dove il terreno digradava verso ovest. Jason ci mostrò la luce zodiacale
che saliva in cielo.»

«Cos'è la luce zodiacale?»

«Luce solare riflessa dai granelli di ghiaccio sulla cintura
asteroidale. A volte si riesce a vederla in una notte buia molto
limpida.» O almeno si riusciva, prima dello Spin. C'era ancora una luce
zodiacale o la pressione solare aveva spazzato via il ghiaccio? «Venne
su dall'orizzonte come il respiro d'inverno, lontanissima e lieve. Diane
era affascinata. Ascoltava Jase che spiegava, a quel tempo le
spiegazioni di Jase la affascinavano ancora - non era ancora cresciuta.
Amava la sua intelligenza, lo amava \emph{per} la sua
intelligenza\ldots»

«Come il padre di Jason, forse? Ora girati sul ventre, per favore.»

«Ma non come se fosse una sua proprietà. Era puro incanto da far
strabuzzare gli occhi.»

«Scusami, ``strabuzzare gli occhi''?»

«Aprire molto gli occhi. Poi il vento cominciò a sollevarsi e Jason
accese la torcia elettrica e la puntò sui pioppi, in modo che Diane
potesse vedere come si muovevano i rami.» Insieme a questa immagine mi
sovvenne un vivido ricordo di Diane giovane con un maglione di almeno
una taglia più grande della sua misura, le mani perse nella lana
lavorata a maglia, che si abbracciava con il viso rivolto all'insù,
verso un cono di luce che le si rifletteva negli occhi, illuminati come
due lune solenni. «Le mostrò come i rami più grandi sbattevano come se
fossero al rallentatore e quelli più piccoli più rapidamente. Questo
succedeva perché ogni ramo e fuscello aveva quella che Jase chiamava una
frequenza di risonanza. E potevi pensare a quelle frequenze di risonanza
come se fossero note musicali, disse. Il movimento dell'albero nel vento
era un tipo di musica dal tono troppo basso per le orecchie umane. Il
tronco cantava una nota grave, i rami recitavano le battute dei tenori e
i fuscelli suonavano l'ottavino. Oppure, disse ancora, si potevano
pensare come numeri puri; ogni risonanza, dal vento al tremore di una
foglia, elabora un calcolo dentro un calcolo dentro un calcolo.»

«Lo stai descrivendo in modo bellissimo» disse Ina.

«Non è bello nemmeno la metà di come lo aveva raccontato Jason. Era come
se fosse innamorato del mondo, o almeno dei suoi schemi\ldots{} della
musica che porta dentro. Ahia!»

«Scusami. E Diane amava Jason?»

La osservai perplesso. «Amava essere sua sorella. Era fiera di lui.»

«E tu amavi essere suo amico?»

«Suppongo di sì.»

«E amavi Diane.»

«Sì.»

«E lei amava te.»

«Forse. Io ci speravo.»

«Poi, se posso chiedere, cosa è andato storto?»

«Cosa ti fa pensare che qualcosa sia andato storto?»

«È ovvio che siete ancora innamorati. Voi due intendo. Ma non come un
uomo e una donna che sono stati insieme per molti anni. Qualcosa deve
avervi tenuti lontani. Perdonami, è davvero impertinente da parte mia.»

Sì, qualcosa ci aveva tenuti lontani. Molte cose. Chiaramente era stato
lo Spin, o almeno era ciò che supposi. L'aveva spaventata per motivi che
non ho mai capito del tutto; come se lo Spin fosse una sfida e un
castigo per tutto ciò che l'aveva fatta sentire al sicuro. Ma cos'è che
la faceva sentire al sicuro? L'ordinato avanzare della vita; gli amici,
la famiglia, il lavoro - una sorta di fondamentale partecipazione
emotiva alle cose, che nella Grande Casa di E.D. e Carol Lawson doveva
essere sembrata una fragilità, più desiderata che reale.

La Grande Casa l'aveva tradita e alla fine anche Jason l'aveva tradita.
Le idee scientifiche che le presentava come doni singolari e che una
volta le erano parse rassicuranti - gli accoglienti accordi maggiori di
Newton ed Euclide - divennero estranee e sempre più alienanti: la
lunghezza di Planck (sotto la quale le cose non si comportavano più come
\emph{cose}), i buchi neri, sigillati dalla loro stessa imponderabile
densità in un regno oltre la causa e l'effetto, un universo che non solo
si espandeva, ma accelerava verso il suo stesso decadimento. Una volta
mi disse, quando Sant'Agostino era ancora vivo, che quando metteva le
mani sulla pelliccia del cane voleva sentirne il calore e la vivacità -
non contare i battiti del cuore o prendere in considerazione gli spazi
tra i nuclei e gli elettroni che costituivano il suo essere fisico.
Voleva che San Cane fosse se stesso e intero, non la somma delle sue
spaventose parti, non un fugace epifenomeno evolutivo nella vita di una
stella morente. C'erano già troppo poco amore e affetto nella sua vita e
ogni esempio di quei sentimenti doveva essere conteggiato e riposto in
paradiso, accumulato per combattere l'inverno dell'universo.

Lo Spin, quando arrivò, doveva esserle sembrato una vendetta mostruosa
della visione del mondo di Jason - ancora di più perché ne era
ossessionato. Chiaramente c'era vita intelligente in altre parti della
galassia; e, in modo altrettanto ovvio, non assomigliava per niente alla
nostra. Era immensamente potente, spaventosamente paziente e del tutto
indifferente al terrore che aveva inflitto al mondo. Immaginando gli
Ipotetici, era possibile figurarsi robot intelligenti o esseri di
energia imperscrutabile; ma mai il tocco di una mano, un bacio, un letto
caldo o una parola di conforto.

Perciò aveva odiato lo Spin in modo profondamente personale e penso che
in sostanza quell'odio l'abbia spinta verso Simon Townsend e il
movimento NR. Nella teologia dell'NR lo Spin divenne un evento sacro ma
anche subordinato: grande ma non grande quanto il Dio di Abramo;
scioccante ma meno sbalorditivo di un Salvatore crocifisso e di una
tomba vuota.

Dissi un po' di queste cose a Ina. Lei rispose: «Ovviamente non sono
cristiana. Non sono neppure islamica a sufficienza da soddisfare le
autorità locali. Corrotta dall'occidente ateo, eccomi qua. Ma persino
nell'Islam ci sono movimenti del genere. Gente che parla a vanvera
dell'Imam Mehdi e Ad-Dajjal, Yajuj e Majuj che bevono il mare della
Galilea. Perché pensavano che tutto questo avesse più senso. Ecco. Ho
finito.» Mi aveva strofinato la pianta dei piedi. «Hai sempre saputo
queste cose di Diane?»

``Saputo in che senso?'' mi chiesi. Percepito, sospettato, intuito; ma
\emph{saputo}\ldots{} no, non potevo dire che fosse così.

«Allora forse il farmaco marziano è all'altezza delle tue aspettative»
disse Ina mentre usciva con la sua bacinella d'acciaio inossidabile
piena d'acqua calda e il suo assortimento di spugne, lasciandomi
qualcosa a cui pensare nella notte buia.

\separatorefiori

C'erano tre porte che conducevano dentro o fuori dalla clinica medica di
Ibu Ina. Una volta mi aveva fatto fare un giro dell'edificio, dopo che
l'ultimo paziente previsto era andato via con un dito steccato.

«Questo è ciò che ho costruito nella mia vita» aveva detto. «Un po'
poco, penserai. Ma la gente di questo villaggio aveva bisogno di
qualcosa tra qui e l'ospedale di Padang\ldots{} una distanza
considerevole se devi viaggiare in autobus oppure se le strade sono
inaffidabili.»

Una porta era quella principale da cui i pazienti andavano e venivano.

Un'altra era sul retro, rivestita di metallo e solida. Ina parcheggiava
la sua piccola automobile elettrica nel parcheggio di terra battuta
dietro la clinica e usava quella porta quando arrivava la mattina e la
chiudeva uscendo alla sera. Era adiacente alla stanza in cui vivevo e
avevo imparato a riconoscere il suono delle sue chiavi che tintinnavano
non molto dopo la prima chiamata alla preghiera della moschea del
villaggio che sarà stata distante poche centinaia di metri.

La terza porta era laterale, in fondo a un piccolo corridoio che
ospitava anche il bagno e una fila di dispense per le forniture. Era là
che accoglieva le consegne e quello era il tragitto che En preferiva per
entrare e uscire.

En era proprio come lo aveva descritto Ina: schivo ma brillante e
abbastanza intelligente da guadagnarsi la laurea in medicina su cui
aveva riposto tutte le sue speranze più intime. I suoi genitori non
erano ricchi, diceva Ina, ma se avesse ottenuto una borsa di studio,
fatto i corsi propedeutici presso la nuova università di Padang,
primeggiato, trovato un modo per farsi finanziare una specializzazione
post-laurea\ldots{} «Poi, chissà? Il villaggio potrebbe avere un nuovo
medico. È così che ho fatto io.»

«Credi che tornerebbe indietro per praticare qui?»

«Potrebbe. Andiamo via, torniamo indietro.» Aveva scrollato le spalle
come se questo fosse l'ordine naturale delle cose. E per i Minangkabau
lo era. \emph{Rantau}, la tradizione di mandare lontano gli uomini
giovani, faceva parte del sistema dell'\emph{adat}, tradizione e regola
ferrea. L'\emph{adat}, come l'Islam conservatore, era stato eroso dagli
ultimi trent'anni di modernità, ma pulsava sotto la superficie della
vita dei Minangkabau come un battito cardiaco.

En era stato diffidato dal disturbarmi ma a poco a poco perse il timore
che provava nei miei confronti. Con l'autorizzazione esplicita di Ibu
Ina, tra un attacco di febbre e l'altro, En affinava il suo vocabolario
inglese portandomi generi alimentari e dicendone il nome per me:
\emph{silomak}, riso appiccicoso; \emph{singgang ayam}, pollo al curry.
Quando dicevo «Grazie», En gridava «Prego!» e ridacchiando mostrava una
serie di denti bianchi e splendenti ma molto irregolari: Ina stava
cercando di convincere i suoi genitori a fargli mettere l'apparecchio.

Anche Ina viveva in una piccola casa nel villaggio con dei parenti,
sebbene ultimamente dormisse in un ambulatorio della clinica, uno spazio
che non doveva essere molto più confortevole della mia triste cella.
Tuttavia alcune sere veniva richiamata a svolgere faccende di famiglia;
quando capitava, annotava la mia temperatura e le mie condizioni, mi
riforniva di cibo e acqua e mi lasciava un cercapersone in caso
d'emergenza. E rimanevo da solo fino al tintinnio della sua chiave nella
porta la mattina dopo.

Ma una notte mi risvegliai dopo un sogno folle e labirintico al suono
della porta laterale che cigolava come se qualcuno avesse girato il
pomello e stesse tentando di aprirla. Non era Ina. Porta sbagliata, ora
sbagliata. Il mio orologio segnava mezzanotte. Stava iniziando la parte
più profonda della notte e in giro c'erano solo pochi abitanti del
villaggio, che infestavano i locali \emph{warung}, qualche macchina che
transitava sulla strada principale e i camion che cercavano di
raggiungere qualche \emph{desa} lontana prima del mattino. Forse era un
paziente che sperava di trovare Ina ancora qui. O forse un drogato a
caccia di dosi.

Il cigolio si fermò.

Mi alzai senza fare rumore e indossai i jeans e una maglietta. La
clinica era al buio, così come la mia cella, e l'unica luce era quella
lunare che penetrava dalla finestra alta\ldots{} che venne
improvvisamente eclissata.

Guardai in alto e vidi la sagoma della testa di En come un pianeta
sospeso. «Pak Tyler!» sussurrò.

«En! Mi hai spaventato.» Infatti lo shock mi aveva prosciugato le forze
dalle gambe. Dovetti appoggiarmi alla parete per stare in piedi.

«Fammi entrare!» disse En.

Quindi camminai silenziosamente a piedi nudi fino alla porta laterale e
aprii il chiavistello.

La brezza che entrò era calda e umida. En entrò subito dopo. «Fammi
parlare con Ibu Ina!»

«Non c'è. Cos'è successo, En?»

Era profondamente sconcertato. Spinse gli occhiali fin sopra la gobbetta
sul naso. «Ma io devo parlare con lei!»

«Stanotte è a casa sua. Sai dove abita?»

En annuì deluso. «Ma mi aveva detto di venire \emph{qui} e parlare con
lei.»

«Di cosa? Ma quando te l'ha detto?»

«Se uno straniero chiede della clinica, devo venire qui a riferirlo a
lei.»

«Ma lei non è\ldots» Poi l'importanza delle sue parole fece breccia
nella nebbia della febbre incipiente. «En, c'è qualcuno in città che fa
domande su Ibu Ina?»

Lo persuasi a raccontarmi la storia. En viveva con la famiglia in una
casa dietro un \emph{warung} (un chiosco del cibo) nel cuore del
villaggio, a una distanza di tre porte dall'ufficio del sindaco, il
\emph{kepala desa}. En, nelle notti insonni, sdraiato in camera sua
riusciva a sentire il mormorio delle conversazioni dei clienti del
\emph{warung}. Così aveva acquisito un archivio enciclopedico per quanto
confuso dei pettegolezzi del villaggio. A tarda sera solitamente erano
gli uomini che sedevano a chiacchierare e a bere caffè, il padre di En,
gli zii e i vicini. Ma quella sera c'erano due stranieri arrivati con
un'elegante macchina nera che si erano avvicinati alle luci del
\emph{warung}, sfrontati come bufali d'acqua, e avevano chiesto, senza
prima presentarsi, dove fosse la clinica locale. Nessuno dei due era
malato. Indossavano abiti di città, si comportavano in modo maleducato e
si muovevano come poliziotti, e così le indicazioni che ricevettero dal
padre di En furono vaghe e sbagliate, e li avrebbero condotti
esattamente da tutt'altra parte.

Ma stavano cercando la clinica di Ina e inevitabilmente l'avrebbero
trovata; in un villaggio questo genere d'informazioni errate poteva al
massimo provocare un ritardo. Allora En si era vestito e, senza farsi
notare, era uscito di corsa da casa sua per raggiungere la clinica,
seguendo le istruzioni ricevute e rispettando il patto stretto con Ibu
Ina di avvertirla in caso di pericolo.

«Bravissimo» gli dissi. «Ottimo lavoro, En. Ora devi andare a casa sua e
raccontarle tutto.»

E nel frattempo avrei raccolto le mie cose e sarei uscito dalla clinica.
Immaginavo di potermi nascondere nelle risaie adiacenti finché la
polizia non fosse arrivata e andata via. Ero abbastanza forte per farlo.
Forse.

Ma En incrociò le braccia e si allontanò da me. «Mi ha detto di
aspettarla \emph{qui}.»

«Va bene, ma non arriverà prima di domattina.»

«Dorme qui quasi tutte le notti.» Allungò il collo guardando dietro di
me lungo il corridoio buio della clinica, come se lei stesse per uscire
dall'ambulatorio per rassicurarlo.

«Sì, ma non stanotte. Giuro, En, potrebbe essere pericoloso. Queste
persone potrebbero essere nemiche di Ibu Ina, capisci?»

Era dominato da una certa testardaggine innata e fiera. Per quanto
fossimo diventati amici, En non si fidava ancora di me. Tremò per un
attimo, gli occhi spalancati come un lemure, poi mi aggirò sfrecciando e
si addentrò nella clinica illuminata dalla Luna, chiamando: «Ina! Ina!»

Lo rincorsi e spostandomi accendevo le luci.

Allo stesso tempo cercavo di pensare lucidamente alla situazione. Gli
uomini maleducati che cercavano la clinica potevano far parte della
nuova Reformasi di Padang, oppure essere poliziotti locali, o forse
lavoravano per l'Interpol o per il Dipartimento di Stato o per qualsiasi
altra agenzia scelta dall'amministrazione Chaykin per piombarci addosso
come un martello.

E se erano lì per cercare me, significava forse che avevano trovato e
interrogato Jala, l'ex marito di Ina? Voleva dire che avevano già
arrestato Diane?

En si ritrovò dentro un ambulatorio buio. Batté la fronte sulle staffe
aperte di un lettino e finì col sedere per terra. Quando lo raggiunsi
stava piangendo senza far rumore, spaventato, e le lacrime gli colavano
sulle guance. Il segno sull'occhio sinistro aveva un brutto aspetto ma
non era pericoloso.

Gli poggiai le mani sulle spalle. «En, Ina non è qui, sul serio. Devi
credermi, non c'è, davvero. E sono sicuro che non intendeva dirti di
correre qui di notte col rischio che ti capitasse qualcosa di brutto.
Lei non lo permetterebbe, o sbaglio?»

«Uh» borbottò En, dandomi ragione.

«Allora corri a casa, va bene? Corri a casa e resta là. Mi occuperò io
di questo problema e domani vedremo insieme Ibu Ina. Ti sembra
ragionevole?»

En provò a sostituire la sua paura con un approccio critico. «Credo di
sì» disse storcendo il naso.

Lo aiutai ad alzarsi.

Ma poi ci fu un suono di ghiaia che scricchiolava sotto le ruote davanti
alla clinica e ci accucciammo di nuovo entrambi.

\separatorefiori

Ci precipitammo nella stanza dell'accettazione. Sbirciai fuori tra i
listelli di bambù delle tendine. En stava dietro di me, le piccole mani
annodate alla stoffa della mia camicia.

La macchina rimase col motore al minimo al chiaro di Luna. Non riconobbi
il modello ma a giudicare dal nero brillante sembrava piuttosto nuova.
Ci fu un breve bagliore nel buio all'interno, come di un accendino. Poi
una luce molto più intensa, un faretto abbagliante puntato dal
finestrino passeggero. Passò attraverso le tendine e proiettò delle
ombre ondulate sui manifesti con le norme igieniche sulla parete
opposta. Abbassammo la testa. En piagnucolava.

«Pak Tyler?» disse.

Chiusi gli occhi e scoprii che era difficile riaprirli. Dietro le
palpebre vidi delle girandole e delle stelline scoppiettanti. Ancora la
febbre. Un piccolo coro di vocine interne ripeteva: ``Ancora la febbre,
ancora la febbre'' prendendomi in giro.

«Pak Tyler!» ripeté En.

Davvero un pessimo tempismo. (``Pessimo tempismo, pessimo
tempismo\ldots'')

«Vai alla porta, En. La porta laterale.»

«Vieni con me!»

Ottimo consiglio. Controllai nuovamente la finestra. Il faretto si era
spento. Mi alzai e condussi En lungo il corridoio e oltre gli armadietti
dei medicinali fino alla porta laterale, che aveva lasciato aperta. La
sera era ingannevolmente tranquilla, ingannevolmente invitante; una
spanna di terra battuta, una risaia; la foresta, palme che sembravano
nere al chiaro di Luna e che scuotevano le chiome dolcemente.

L'edificio della clinica stava tra noi e la macchina. «Corri dritto
nella foresta» dissi.

«Conosco la via\ldots»

«Stai lontano dalla strada. Nasconditi se devi.»

«Lo so. Vieni con me!»

«Non posso» dissi, letteralmente. Nelle mie attuali condizioni l'idea di
correre dietro a un ragazzino di dieci anni era assurda.

«Ma\ldots» disse En, e gli detti una spinta leggera esortandolo a non
perdere tempo.

Corse via senza voltarsi sparendo nell'ombra a una velocità quasi
inquietante, silenzioso, piccolo, ammirevole. Lo invidiavo. Nella quiete
che seguì, sentii lo sportello di una macchina che si apriva e si
richiudeva.

La Luna era piena per tre quarti, più rubiconda e distante che mai, con
una faccia diversa da quella dei miei ricordi d'infanzia. Niente più
Uomo nella Luna, il gioco che facevamo da piccoli, immaginando che le
macchie lunari formassero il disegno di un uomo; e quella scura
cicatrice ovoidale lungo la sua superficie, quel nuovo ma ora antico
\emph{mare}, era il risultato di un impatto massiccio che aveva sciolto
la regolite dal polo all'equatore e aveva rallentato la graduale spirale
della Luna in allontanamento dalla Terra.

Dietro di me, sentii i poliziotti bussare alla porta d'ingresso e
annunciarsi in modo burbero, scuotendo il pomello.

Pensai di correre. Contavo di poter correre - non abilmente come En, ma
con risultati positivi - almeno fino alla risaia. E nascondermi là e
sperare in bene.

Ma poi pensai ai bagagli che avevo lasciato nella stanza sul retro di
Ina. Le valigie non contenevano solo vestiti, ma taccuini e dischetti,
piccoli nastri di memoria digitale e incriminanti fialette di liquido
chiaro.

Tornai dentro e mi richiusi la porta alle spalle. Camminavo scalzo e
all'erta, cercando di captare i rumori dei poliziotti. Forse stavano
girando attorno all'edificio oppure tentavano di entrare dalla porta
d'ingresso. Sentivo che la febbre stava salendo in fretta, non potevo
sapere con certezza se tutti i suoni che udivo erano reali.

Nella stanza nascosta di Ina la luce sul soffitto era ancora fuori uso.
Mi mossi a tastoni sfruttando la luce della Luna. Aprii una delle due
valigette dal corpo rigido e ci infilai una pila di pagine scritte a
mano; la chiusi a chiave, la sollevai e barcollai. Poi presi la seconda
come zavorra di contrappeso e scoprii che riuscivo a malapena a
camminare.

Stavo quasi per inciampare su un piccolo oggetto di plastica che poi
riconobbi essere il cercapersone di Ina. Mi fermai, misi giù le valigie,
afferrai il cercapersone e lo infilai nella tasca della camicia. Poi
tirai qualche respiro profondo e di nuovo sollevai le valigie;
misteriosamente, sembravano essere diventate più pesanti. Cercai di
ripetere a me stesso ``ce la puoi fare'', ma quelle parole erano banali
e poco convincenti e riecheggiarono nel mio cranio come se fosse
diventato grande quanto una cattedrale.

Sentii dei rumori alla porta posteriore, quella che Ina teneva chiusa
con un lucchetto esterno: il clangore del metallo e il gemito del
chiavistello, forse un piede di porco inserito tra le cerniere del
lucchetto e fatto girare. E molto presto, inevitabilmente, il lucchetto
avrebbe ceduto e gli uomini di quella macchina sarebbero entrati.

Barcollai verso la terza porta, quella laterale che usava En, la
sbloccai e la aprii sperando ciecamente che non ci fosse nessuno ad
aspettarmi fuori. Non c'era nessuno. Entrambi gli intrusi (sempre che
fossero solo due) erano sul retro. Bisbigliavano, mentre armeggiavano
con il lucchetto. Le loro voci si percepivano appena sopra al coro di
rane e al lieve suono del vento.

Non ero certo di riuscire a nascondermi nella risaia senza essere visto.
Peggio ancora, non ero certo di riuscirci senza cadere a terra, ma
quando il lucchetto si staccò dalla porta ci fu un forte schianto. ``Lo
sparo dello starter'' pensai. ``Ce la puoi fare'' pensai. Raccolsi le
valigie e barcollai scalzo sotto la notte stellata.